MongoDB to otwartoźródłowy dokumentowy system zarządzania bazą danych
rozwijany od 2007 roku przez firmę MongoDB Inc. (dawniej 10gen). Obecna
linia 7.x (stan na 2025~r.) należy do kategorii baz NoSQL i zamiast
tabel z wierszami przechowuje dane w postaci dokumentów BSON (Binary JSON),
które mogą mieć dowolną, zagnieżdżoną strukturę. MongoDB jest dostępny
na licencji SSPL (Server Side Public License) dla wersji Community oraz
komercyjnej dla wersji Enterprise.

\subsection{Skalowalność}

MongoDB zostało zaprojektowane z myślą o skalowalności poziomej jako
natywnej cesze silnika. Wbudowany mechanizm shardingu pozwala automatycznie
rozdzielać dane między wiele węzłów na podstawie wybranego klucza podziału
(shard key), co umożliwia obsługę zbiorów danych przekraczających pojemność
pojedynczego serwera. Warstwa routingu (mongos) jest transparentna dla
aplikacji -- zapytania są kierowane do właściwych węzłów automatycznie.

Replikacja realizowana jest przez Replica Set, czyli grupę węzłów składającą
się z jednego primary i co najmniej dwóch secondary. Wszystkie zapisy trafiają
do primary, skąd są asynchronicznie replikowane do secondary. Odczyty mogą
być opcjonalnie kierowane do secondary, co odciąża primary przy
obciążeniach zdominowanych przez odczyty. Skalowanie pionowe również jest
obsługiwane -- MongoDB efektywnie wykorzystuje pamięć RAM przez agresywne
cachowanie często używanych danych w pamięci operacyjnej.

\subsection{Bezpieczeństwo}

MongoDB oferuje wielowarstwowy model bezpieczeństwa. Uwierzytelnianie
obsługuje hasła (domyślnie SCRAM-SHA-256), certyfikaty X.509 oraz
integrację z LDAP i Kerberos (w Enterprise Edition). Kontrola dostępu
oparta jest na rolach (RBAC) -- wbudowane role obejmują odczyt, zapis,
administrację i inne, a użytkownik może definiować własne role z precyzyjnymi
uprawnieniami do kolekcji i operacji.

Komunikacja między klientem a serwerem oraz między węzłami klastra może
być szyfrowana przez TLS. MongoDB Enterprise oferuje szyfrowanie danych
w spoczynku (Encrypted Storage Engine) oraz zaawansowany audyt operacji.
Od wersji 4.2 dostępne jest też szyfrowanie na poziomie pola (Field Level
Encryption), które pozwala szyfrować wybrane pola dokumentu po stronie
klienta jeszcze przed wysłaniem do serwera -- serwer nigdy nie widzi
niezaszyfrowanych wartości.

\subsection{Awaryjność}

Podstawą trwałości danych w MongoDB jest dziennik (journal), który działa
podobnie do WAL w silnikach relacyjnych -- każda operacja jest najpierw
zapisywana do dziennika, co gwarantuje możliwość odtworzenia stanu po
awarii. Silnik składowania WiredTiger (domyślny od wersji 3.2) obsługuje
transakcje ACID na poziomie dokumentu natywnie, a od wersji 4.0 również
transakcje wielodokumentowe i wielokolekcyjne.

Replica Set zapewnia automatyczny failover -- jeśli primary przestanie
odpowiadać, węzły secondary przeprowadzają wybory i jeden z nich przejmuje
rolę primary, zazwyczaj w ciągu kilku do kilkunastu sekund. Proces ten
jest w pełni automatyczny i transparentny dla aplikacji korzystających
ze sterowników obsługujących retryable writes. Kopie zapasowe można
tworzyć przez \texttt{mongodump} (logiczne) lub snapshoty systemu plików
(fizyczne, wymagają zamrożenia operacji zapisu).

\subsection{Migracje}

Do eksportu i importu danych służą narzędzia \texttt{mongodump} i
\texttt{mongorestore} (format BSON) oraz \texttt{mongoexport} i
\texttt{mongoimport} (format JSON/CSV, bardziej przenośny między systemami).
Migracja między wersjami MongoDB jest zazwyczaj prosta dzięki rolling
upgrade -- węzły klastra aktualizuje się po kolei bez zatrzymywania całego
systemu.

Migracja z relacyjnych baz danych jest bardziej złożona, bo wymaga
przeprojektowania schematu z modelu tabelarycznego na dokumentowy.
W praktyce często nie jest to proste przeniesienie danych, lecz zmiana
podejścia do modelowania -- na przykład połączenie kilku tabel w jeden
zagnieżdżony dokument. Narzędzia takie jak Mongify czy własne skrypty ETL
wspierają ten proces. Wersjonowanie schematu dokumentów, które z natury
jest elastyczne, realizuje się przez zewnętrzne narzędzia lub własną logikę
walidacji w aplikacji (np. przez JSON Schema Validation dostępne w MongoDB).

\subsection{Integracje}

MongoDB oferuje oficjalne sterowniki dla wszystkich popularnych języków
programowania: Python (pymongo, motor), Node.js (mongodb, mongoose),
Java, Go, C\#/.NET, Ruby, PHP i innych. Mongoose (Node.js) i MongoEngine
(Python) są najpopularniejszymi bibliotekami ODM (Object Document Mapper)
ułatwiającymi pracę z dokumentami.

Narzędzia analityczne integrują się z MongoDB przez konektor BI (MongoDB
Connector for BI), który tłumaczy zapytania SQL na operacje MongoDB,
co pozwala na użycie narzędzi jak Tableau czy Power BI bez modyfikacji.
MongoDB Atlas (chmurowa wersja zarządzana) jest dostępny na AWS, Google
Cloud i Azure i oferuje dodatkowe funkcje jak Atlas Search (pełnotekstowe
wyszukiwanie oparte na Apache Lucene), Atlas Vector Search (wyszukiwanie
wektorowe dla zastosowań AI) oraz Charts (wbudowana wizualizacja danych).

\subsection{Obszary biznesowe}

MongoDB najlepiej sprawdza się w scenariuszach, gdzie dane mają zmienną
lub hierarchiczną strukturę i gdzie elastyczność schematu jest istotna:

\begin{itemize}
    \item \textbf{Platformy treści i katalogi produktów} -- dokumentowy
    model danych pozwala przechowywać produkty o różnych atrybutach
    w jednej kolekcji bez konieczności tworzenia dziesiątek tabel lub
    kolumn nullable. Stosowany przez serwisy takie jak eBay czy Etsy.
    \item \textbf{Aplikacje mobilne i IoT} -- elastyczny schemat i
    natywna obsługa JSON ułatwiają przechowywanie danych o zmiennej
    strukturze napływających z urządzeń mobilnych i sensorów.
    \item \textbf{Systemy zarządzania treścią (CMS)} -- zagnieżdżone
    dokumenty dobrze odwzorowują hierarchiczną strukturę artykułów,
    komentarzy i metadanych w jednym obiekcie.
    \item \textbf{Analityka w czasie rzeczywistym} -- agregacje MongoDB
    i możliwość przechowywania surowych danych zdarzeń bez narzutu
    normalizacji przyspieszają budowę pipeline'ów analitycznych.
    \item \textbf{Personalizacja i profile użytkowników} -- zmienna
    struktura dokumentu pozwala przechowywać preferencje, historię
    i ustawienia różnych użytkowników bez modyfikacji schematu.
\end{itemize}

\subsection{Zalety}

\begin{itemize}
    \item \textbf{Elastyczny schemat} -- dokumenty w kolekcji mogą mieć
    różną strukturę. Dodanie nowego pola nie wymaga migracji schematu,
    co przyspiesza iteracje w fazie rozwoju aplikacji.
    \item \textbf{Natywny model dokumentowy} -- zagnieżdżone dokumenty
    i tablice pozwalają przechowywać powiązane dane razem, eliminując
    potrzebę kosztownych JOIN-ów przy odczycie. Jeden dokument to często
    kompletny obiekt biznesowy.
    \item \textbf{Natywny sharding} -- skalowanie poziome jest wbudowaną
    cechą silnika, a nie rozwiązaniem zewnętrznym jak w PostgreSQL czy MySQL.
    \item \textbf{Wysoka wydajność zapisu} -- brak narzutu normalizacji
    i możliwość aktualizacji zagnieżdżonych pól przez operatory
    (\texttt{\$set}, \texttt{\$push}, \texttt{\$inc}) bez przepisywania
    całego dokumentu przekłada się na dobrą wydajność przy intensywnym
    zapisie.
    \item \textbf{Bogaty język zapytań} -- agregacje (Aggregation Pipeline)
    pozwalają na złożone transformacje danych bezpośrednio w bazie,
    a Atlas Search dodaje pełnotekstowe wyszukiwanie klasy Elasticsearch.
    \item \textbf{Atlas i ekosystem chmurowy} -- MongoDB Atlas oferuje
    w pełni zarządzany klaster z automatycznym skalowaniem, backupami
    i monitoringiem, dostępny na trzech głównych chmurach.
\end{itemize}

\subsection{Wady i ograniczenia}

\begin{itemize}
    \item \textbf{Brak JOIN-ów na poziomie silnika} -- relacje między
    kolekcjami realizuje się przez \texttt{\$lookup} (odpowiednik LEFT JOIN),
    który jest znacznie wolniejszy od JOIN-ów w silnikach relacyjnych,
    szczególnie przy złożonych powiązaniach. Model dokumentowy wymaga
    świadomego projektowania pod kątem wzorców dostępu.
    \item \textbf{Ograniczone transakcje wielodokumentowe} -- choć dostępne
    od wersji 4.0, transakcje obejmujące wiele dokumentów i kolekcji
    mają znacznie wyższy narzut niż w silnikach relacyjnych i są odradzane
    jako główny wzorzec użycia.
    \item \textbf{Brak schematu jako wada} -- elastyczność schematu może
    być pułapką. Brak wymuszanej struktury prowadzi do niespójności danych,
    jeśli walidacja nie jest implementowana na poziomie aplikacji lub przez
    JSON Schema Validation.
    \item \textbf{Wyższe zużycie pamięci} -- BSON przechowuje nazwy pól
    w każdym dokumencie, co przy kolekcjach z wieloma polami zwiększa
    rozmiar danych w porównaniu do tabel relacyjnych.
    \item \textbf{Licencja SSPL} -- zmiana licencji z Apache 2.0 na SSPL
    w 2018 roku jest kontrowersyjna i dla części organizacji wyklucza
    użycie MongoDB Community Edition w pewnych scenariuszach.
    \item \textbf{Słabsze narzędzia analityczne} -- MongoDB nie jest
    silnikiem OLAP i przy złożonych zapytaniach analitycznych na dużych
    zbiorach ustępuje specjalizowanym rozwiązaniom jak ClickHouse czy
    BigQuery.
\end{itemize}

\subsection{Udogodnienia}

\begin{itemize}
    \item \textbf{MongoDB Compass} -- oficjalne graficzne narzędzie do
    przeglądania danych, budowania zapytań i analizy planów wykonania
    (Explain Plan). Pozwala wizualnie budować pipeline agregacji bez
    pisania kodu.
    \item \textbf{mongosh} -- nowoczesna powłoka interaktywna oparta
    na Node.js zastępująca starszy klient \texttt{mongo}. Obsługuje
    JavaScript, autouzupełnianie i kolorowanie składni.
    \item \textbf{Aggregation Pipeline} -- jeden z najbardziej elastycznych
    mechanizmów transformacji danych w bazach NoSQL. Pozwala łączyć
    filtrowanie, grupowanie, sortowanie, łączenie kolekcji i obliczenia
    w jednym zapytaniu wykonywanym po stronie serwera.
    \item \textbf{Change Streams} -- mechanizm subskrypcji na zmiany
    w kolekcji lub bazie w czasie rzeczywistym, przydatny do budowy
    reaktywnych aplikacji i pipeline'ów CDC bez zewnętrznych narzędzi.
    \item \textbf{JSON Schema Validation} -- możliwość wymuszenia struktury
    dokumentów na poziomie kolekcji przez zdefiniowanie schematu JSON Schema,
    co pozwala połączyć elastyczność dokumentową z walidacją danych.
    \item \textbf{Atlas Search i Vector Search} -- wbudowane wyszukiwanie
    pełnotekstowe klasy Elasticsearch oraz wyszukiwanie wektorowe dla
    zastosowań AI/ML, dostępne bez konieczności integracji zewnętrznych
    systemów.
\end{itemize}